\subsection{Проектирование интерфейсов взаимодействия и сценариев работы системы}

Интерфейсы проектируемой системы строятся вокруг трех основных контуров
взаимодействия: контура получения данных от узла Bitcoin, контура
административного управления и контура прикладной выдачи данных внешним
потребителям. Такое разделение позволяет формализовать назначение каждой группы
операций и уменьшить связанность между внешними участниками системы.

Контур получения данных ориентирован на доверенное взаимодействие с Bitcoin
Core через RPC-прокси. Он предназначен исключительно для внутреннего
использования индексером и не должен быть доступен прикладным клиентам. Контур
административного управления реализуется через REST API, административную
панель и CLI. В его рамках выполняются операции запуска, остановки, паузы и
возобновления заданий индексирования, а также получение сведений о состоянии узлов и
выполняемых процессов. Контур прикладной выдачи предоставляет доступ к блокам,
транзакциям, балансам адресов, UTXO и mempool-данным.

Основные проектируемые группы REST endpoints включают:
\begin{enumerate}
    \item endpoints управления заданиями индексирования;
    \item endpoints мониторинга и диагностики узлов;
    \item endpoints получения блоков, транзакций, балансов и UTXO;
    \item сервисные endpoints наблюдаемости и документирования API.
\end{enumerate}

В таблице~\ref{tab:api-groups} приведены основные группы прикладных endpoint'ов
и их назначение.

\begin{footnotesize}
\setlength{\tabcolsep}{3pt}
\begin{longtable}{|p{0.24\textwidth}|p{0.34\textwidth}|p{0.28\textwidth}|}
    \caption{Группы REST API системы}\label{tab:api-groups}\\ \hline
    Группа & Endpoint'ы & Назначение \\ \hline
    \endfirsthead
    \caption*{Продолжение таблицы \ref{tab:api-groups}}\\ \hline
    Группа & Endpoint'ы & Назначение \\ \hline
    \endhead
    \hline
    \endfoot
    \hline
    \endlastfoot
    Jobs API &
    \texttt{/v1/jobs}, \texttt{/v1/jobs/\{id\}}, операции \texttt{start}, \texttt{stop}, \texttt{pause}, \texttt{resume}, \texttt{retry} &
    создание, просмотр и управление заданиями индексирования \\ \hline
    Nodes API &
    \texttt{/v1/nodes}, \texttt{/v1/nodes/\{id\}} &
    получение сведений о подключенных узлах и их состоянии \\ \hline
    Data API &
    \texttt{/v1/data/addresses/*}, \texttt{/v1/data/transactions}, \texttt{/v1/data/blocks} &
    прикладная выдача балансов, UTXO, транзакций и блоков \\ \hline
    Observability &
    \texttt{/metrics}, OpenAPI-документация &
    мониторинг и машинно-читаемое описание контрактов API \\ \hline
\end{longtable}
\end{footnotesize}

Жизненный цикл задания индексирования формализуется набором допустимых
переходов, приведенных в таблице~\ref{tab:job-lifecycle}. Такая фиксация важна
для единообразного поведения REST API, CLI и административной панели.

\begin{footnotesize}
\setlength{\tabcolsep}{3pt}
\begin{longtable}{|p{0.17\textwidth}|p{0.27\textwidth}|p{0.22\textwidth}|p{0.22\textwidth}|}
    \caption{Допустимые переходы состояний job}\label{tab:job-lifecycle}\\ \hline
    Операция & Переход & Условие & Результат \\ \hline
    \endfirsthead
    \caption*{Продолжение таблицы \ref{tab:job-lifecycle}}\\ \hline
    Операция & Переход & Условие & Результат \\ \hline
    \endhead
    \hline
    \endfoot
    \hline
    \endlastfoot
    \texttt{start} & \texttt{created -> running} & job создан и разрешен к запуску & runner начинает обработку высот \\ \hline
    \texttt{pause} & \texttt{running -> paused} & job находится в работе & прогресс сохраняется, обработка приостанавливается \\ \hline
    \texttt{resume} & \texttt{paused -> running} & job был поставлен на паузу & обработка продолжается с сохраненной высоты \\ \hline
    \texttt{stop} & \shortstack[l]{\texttt{running|paused|}\\\texttt{failed -> created}} & требуется сброс активного выполнения & job возвращается в исходное состояние \\ \hline
    \texttt{retry} & \texttt{failed -> running} & ошибка устранена оператором & job запускается повторно с сохраненным прогрессом \\ \hline
\end{longtable}
\end{footnotesize}

Проектирование интерфейсов опирается на единый набор принципов:
\begin{enumerate}
    \item предоставление данных в прикладной форме, а не в сыром RPC-представлении;
    \item детерминированные форматы ответов и ошибок;
    \item разграничение подтвержденных данных и mempool-представлений;
    \item защита административных функций средствами аутентификации;
    \item пригодность интерфейсов для использования как человеком, так и внешними системами.
\end{enumerate}

Сценарии функционирования системы включают последовательную индексацию цепи,
обновление адресных агрегатов, обслуживание пользовательских запросов и
административное управление ходом обработки. Тем самым интерфейсы
проектируются не как набор разрозненных конечных точек доступа, а как
формализованное средство реализации основных процессов системы. В дальнейшем
именно эти сценарии становятся основой для реализации программных модулей и
для построения программы испытаний разработанного программного средства.
